\documentclass[aspectratio=169]{beamer}
\usetheme{Madrid}

\usepackage{graphicx}
\usepackage{booktabs}
\usepackage{xcolor}

% Custom premium theme colors (Slate / Dark Cyan)
\definecolor{slate}{RGB}{30, 41, 59}
\definecolor{cyan}{RGB}{14, 116, 144}
\definecolor{lightgray}{RGB}{241, 245, 249}

\setbeamercolor{palette primary}{bg=slate,fg=white}
\setbeamercolor{palette secondary}{bg=cyan,fg=white}
\setbeamercolor{palette tertiary}{bg=slate,fg=white}
\setbeamercolor{structure}{fg=slate}
\setbeamercolor{background canvas}{bg=white}

\title[AQI Prediction \& Monitoring]{Hazara Division AQI Prediction System}
\subtitle{Real-Time Environmental Monitoring \& Forecasting}
\author[Ahmed Ali]{Ahmed Ali}
\institute[10Pearls]{10Pearls Internship Project}
\date{June 2026}

\begin{document}

\begin{frame}
  \titlepage
\end{frame}

\begin{frame}{Project Overview}
  \begin{itemize}
    \item \textbf{Objective}: Provide localized hourly tracking and 72-hour air quality forecasts for Hazara Division, KPK.
    \item \textbf{Scope}: Covers 6 key districts: Abbottabad, Haripur, Mansehra, Battagram, Upper Kohistan, and Torghar.
    \item \textbf{Core System}:
      \begin{itemize}
        \item Automated ingestion of weather and air pollutant concentrations.
        \item 6 predictive models (Linear, Ridge, RF, Gradient Boosting, XGBoost, LSTM).
        \item Interactive FastAPI dashboard for real-time visualization and data exports.
      \end{itemize}
  \end{itemize}
\end{frame}

\begin{frame}{Environmental Motivation}
  \begin{block}{The Valley Inversion Phenomenon}
    Unlike open coastal regions, Hazara Division's geography features deep mountain valleys that trap emissions.
  \end{block}
  \begin{itemize}
    \item \textbf{Temperature Inversions}: Cold air settles in valley basins overnight, trapping pollutants close to the ground.
    \item \textbf{Topographical Barriers}: High mountains block horizontal wind dispersion, creating persistent local pollution.
    \item \textbf{Seasonal Drivers}: Winter wood heating, dust, and agricultural waste burning cause sharp spikes in particulate matter.
  \end{itemize}
\end{frame}

\begin{frame}{Data Pipeline \& Ingestion}
  \begin{itemize}
    \item \textbf{Data Source}: Open-Meteo Air Quality API.
    \item \textbf{Target Metric}: US Air Quality Index (US AQI).
    \item \textbf{Pollutants Monitored}:
      \begin{itemize}
        \item Particulate Matter (PM2.5, PM10).
        \item Gaseous species (CO, NO2, SO2, Ozone).
      \end{itemize}
    \item \textbf{Data Cleaning}:
      \begin{itemize}
        \item Handled extreme outliers by clipping at the 1st and 99th percentiles.
        \item Filled missing intervals using forward-fill (FFill) and backward-fill (BFill).
      \end{itemize}
  \end{itemize}
\end{frame}

\begin{frame}{Feature Engineering}
  We engineer 13 dynamic features to capture persistence and trends:
  \begin{itemize}
    \item \textbf{Lag Features}: 1-hour, 2-hour, and 24-hour lags of AQI to model short-term persistence.
    \item \textbf{Rolling Statistics}: 6-hour and 24-hour moving averages for PM2.5 and overall AQI to track background build-up.
    \item \textbf{Rate of Change}: First-order difference of the AQI value to capture build-up velocity.
    \item \textbf{Diurnal \& Seasonal Indicators}: Time of day, day of the week, and month of the year.
  \end{itemize}
\end{frame}

\begin{frame}{Model Comparisons}
  Models evaluated using chronological 80/20 train/test split:
  \begin{center}
    \begin{tabular}{llccc}
      \toprule
      Model Type & Category & RMSE & MAE & $R^2$ Score \\
      \midrule
      \textbf{Linear Regression} & Statistical & \textbf{4.6796} & \textbf{2.2829} & \textbf{0.8659} \\
      Ridge Regression & Regularized & 4.6797 & 2.2831 & 0.8659 \\
      XGBoost & Boosting & 5.8714 & 3.0906 & 0.7891 \\
      Gradient Boosting & Boosting & 5.9938 & 3.2084 & 0.7802 \\
      Random Forest & Ensemble & 7.4568 & 3.8614 & 0.6593 \\
      LSTM & Neural Network & 10.0665 & 7.9150 & 0.3781 \\
      \bottomrule
    \end{tabular}
  \end{center}
\end{frame}

\begin{frame}{Model Insights}
  \begin{itemize}
    \item \textbf{Linear Performance}: Linear and Ridge regression models achieved the lowest error. Short-term air quality is highly linear (correlation coefficient of lag-1h with target AQI exceeds 0.95).
    \item \textbf{Deep Learning Representation}: LSTM models capture temporal dependencies over the past 24 hours. While standalone metrics are higher, it handles non-linear dynamics effectively.
  \end{itemize}
\end{frame}

\begin{frame}{Explainability (SHAP Values)}
  To verify predictions, we calculated global SHAP importance for the primary model:
  \begin{itemize}
    \item \textbf{1. Lag-1h AQI} (Most critical): Dictates immediate persistence (current air directly affects the next hour).
    \item \textbf{2. 24h Rolling AQI}: Represents background accumulations and slow dispersion patterns.
    \item \textbf{3. PM2.5 24h Average}: Confirms fine particulate matter as the major driving pollutant across all districts.
    \item \textbf{4. Diurnal Hour}: Captures morning/evening rush hour patterns and daily thermal cycles.
  \end{itemize}
\end{frame}

\begin{frame}{Dashboard \& Web Interface}
  \begin{itemize}
    \item \textbf{Backend}: FastAPI for high performance, request routing, and asynchronous API endpoints.
    \item \textbf{Frontend UI}: Clean dark theme built with vanilla HTML/CSS.
    \item \textbf{Key Features}:
      \begin{itemize}
        \item District Selection: Quick switching across all six districts.
        \item Model Selection: Interactive leaderboard showing active metrics.
        \item Interactive Charts: Embedded Plotly visualizations for history and forecasts.
        \item Explainability Panel: Live rendering of SHAP feature importance.
        \item Data Export: Single-click download of full 72-hour forecast CSVs.
      \end{itemize}
  \end{itemize}
\end{frame}

\begin{frame}{Production Automation (MLOps)}
  \begin{itemize}
    \item \textbf{Local Controller}: `run.sh` script to manage execution, status checks, and local manual pipelines.
    \item \textbf{Automated Ingestion}: GitHub Actions workflow runs every hour to fetch raw readings, clean data, and update CSV storages.
    \item \textbf{Automated Retraining}: Daily GitHub Actions workflow retrains all models, computes fresh validation metrics, regenerates SHAP plots, and pushes update files.
    \item \textbf{Automatic Redeployments}: Connected to Render hosting platform for instant deployments on commit.
  \end{itemize}
\end{frame}

\begin{frame}{Conclusion \& Future Outlook}
  \begin{itemize}
    \item Successfully built and deployed a production-grade AQI prediction platform.
    \item Solved key localization issues including valley inversion modeling and local timezone alignment.
    \item \textbf{Future Work}:
      \begin{itemize}
        \item Integrate local physical meteorological models (WRF) for wind vectors.
        \item Incorporate holiday and industrial activity schedules into features.
        \item Expand coverage to additional districts in KPK.
      \end{itemize}
  \end{itemize}
\end{frame}

\end{document}
